%%%%%%%%%%%%%%%%%%%%%%%%%%%%%%%%%%%%%%%%%
% Estructura del libro (diseño + normas APA 7.ª edición, en español)
%
% Basado en "The Legrand Orange Book", versión 2.1.1
% Autor original: Mathias Legrand, con modificaciones de Vel (LaTeXTemplates.com)
% Licencia: CC BY-NC-SA 3.0
%
% Cambios: idioma español, citas y referencias APA 7 (biblatex-apa),
% tablas y figuras en formato APA, márgenes de 2,54 cm, Times 12 pt,
% interlineado 1,5 y paleta azul o verde clara en lugar del naranja.
%
% Normalmente NO necesitas editar este archivo: el color se elige en main.tex
% y los datos del libro en datos.tex.
%%%%%%%%%%%%%%%%%%%%%%%%%%%%%%%%%%%%%%%%%

%----------------------------------------------------------------------------------------
%	PAQUETES BÁSICOS Y PÁGINA
%----------------------------------------------------------------------------------------

% APA: márgenes de 2,54 cm (1 pulgada) por todos los lados, papel carta
\usepackage[letterpaper,margin=2.54cm,headsep=10pt,headheight=16pt]{geometry}

\usepackage{graphicx} % Imágenes
\graphicspath{{imagenes/\tema/}{imagenes/}} % Primero busca en la carpeta del color elegido

\usepackage{lipsum} % Texto de relleno (puedes quitarlo cuando ya no lo uses)

\usepackage{tikz} % Dibujos de la portada y los capítulos
\usetikzlibrary{babel} % Evita conflictos entre TikZ y el español

% Idioma español: "Tabla" en vez de "Cuadro", sin atajos que rompan TikZ
\usepackage[english,spanish,es-tabla,es-noshorthands,es-noquoting,es-nodecimaldot]{babel}
\addto\captionsspanish{%
  \renewcommand{\contentsname}{Tabla de contenido}%
  \renewcommand{\listtablename}{Lista de tablas}%
  \renewcommand{\listfigurename}{Lista de figuras}%
}

\usepackage{enumitem} % Listas
\setlist{nolistsep}

\usepackage{booktabs} % Líneas horizontales de tablas (APA: sin líneas verticales)

\usepackage{zref-totpages} % Número total de páginas para la página legal

%----------------------------------------------------------------------------------------
%	PALETA DE COLORES (se elige con \tema en main.tex)
%----------------------------------------------------------------------------------------

\usepackage{xcolor}
\def\temaverde{verde}
\ifx\tema\temaverde
  \definecolor{principal}{RGB}{61,179,134} % Verde claro
\else
  \definecolor{principal}{RGB}{58,154,217} % Azul claro
\fi
% Versión más oscura del color, para textos (mejor contraste sobre blanco)
\colorlet{acento}{principal!70!black}

%----------------------------------------------------------------------------------------
%	FUENTES E INTERLINEADO (APA: Times New Roman 12 pt)
%----------------------------------------------------------------------------------------

\usepackage{avant} % Avant Garde para títulos
\usepackage{mathptmx} % Times para el texto y las matemáticas
\usepackage{microtype}
\usepackage[utf8]{inputenc}
\usepackage[T1]{fontenc}

\usepackage{setspace}
\onehalfspacing % Interlineado 1,5. Usa \doublespacing si te piden doble espacio

\setlength{\parindent}{1.27cm} % APA: sangría de 1,27 cm en la primera línea
% Texto justificado. Si te piden alineado a la izquierda (APA estricto),
% descomenta las dos líneas siguientes:
%\usepackage{ragged2e}
%\AtBeginDocument{\RaggedRight\setlength{\parindent}{1.27cm}}

%----------------------------------------------------------------------------------------
%	CITAS Y REFERENCIAS APA 7
%----------------------------------------------------------------------------------------

\usepackage[autostyle,spanish=mexican]{csquotes} % Comillas “ ” en lugar de « »
\usepackage[style=apa,backend=biber]{biblatex}
\addbibresource{referencias.bib}
% APA usa "&" entre autores dentro del paréntesis y en la lista de referencias.
% Si tu docente pide "y" en todos los casos, descomenta estas dos líneas:
%\DeclareDelimFormat[parencite]{finalnamedelim}{\addspace\bibstring{and}\space}
%\DeclareDelimFormat[bib,biblist]{finalnamedelim}{\ifthenelse{\value{listcount}>\maxprtauth}{}{\addspace\bibstring{and}\space}}

\usepackage{calc}
\usepackage{makeidx}
% Índice alfabético opcional: descomenta \makeindex y \printindex (al final de main.tex)
%\makeindex

%----------------------------------------------------------------------------------------
%	TABLAS Y FIGURAS EN FORMATO APA
%----------------------------------------------------------------------------------------

% APA: "Tabla 1" en negrita, en su propia línea; el título en cursiva debajo;
% todo encima de la tabla o figura. Numeración corrida (1, 2, 3...), no por capítulo.
\usepackage{caption}
\captionsetup{
  labelsep=newline,
  labelfont=bf,
  textfont=it,
  justification=raggedright,
  singlelinecheck=false,
  font={stretch=1},
  position=top,
  skip=4pt
}
\counterwithout{table}{chapter}
\counterwithout{figure}{chapter}

% Nota debajo de una tabla o figura: \nota{Elaboración propia con datos de ...}
\newcommand{\nota}[1]{\par\vspace{4pt}{\raggedright\small\textit{Nota.}\ #1\par}}

% APA: citas de 40 palabras o más, en bloque con sangría y sin comillas
\renewenvironment{quote}{\list{}{\leftmargin=1.27cm\rightmargin=0pt\topsep=6pt}\item\relax}{\endlist}
\newenvironment{citalarga}{\begin{quote}}{\end{quote}}

% Resumen / Abstract
\newcommand{\palabrasclave}[1]{\par\medskip\hspace*{1.27cm}\textit{Palabras clave}: #1\par}
\newcommand{\keywords}[1]{\par\medskip\hspace*{1.27cm}\textit{Keywords}: #1\par}

% Encabezados APA de nivel 4 y 5 (en línea con el párrafo, terminan en punto)
\newcommand{\nivelcuatro}[1]{\par\indent{\sffamily\bfseries #1.}\ }
\newcommand{\nivelcinco}[1]{\par\indent{\sffamily\bfseries\itshape #1.}\ }

%----------------------------------------------------------------------------------------
%	TABLA DE CONTENIDO
%----------------------------------------------------------------------------------------

\usepackage{titletoc}

\contentsmargin{0cm}
\setcounter{tocdepth}{2} % Muestra capítulos, secciones y subsecciones (1, 1.1, 1.1.1)

% Partes
\titlecontents{part}[0cm]
{\addvspace{20pt}\centering\large\bfseries}
{}
{}
{}

% Capítulos
\titlecontents{chapter}[1.25cm]
{\addvspace{12pt}\large\sffamily\bfseries}
{\color{principal}\contentslabel[\Large\thecontentslabel]{1.25cm}\color{acento}}
{\color{acento}}
{\color{principal}\normalsize\;\titlerule*[.5pc]{.}\;\thecontentspage}

% Secciones
\titlecontents{section}[1.25cm]
{\addvspace{3pt}\sffamily\bfseries}
{\contentslabel[\thecontentslabel]{1.25cm}}
{}
{\hfill\color{black}\thecontentspage}
[]

% Subsecciones
\titlecontents{subsection}[1.25cm]
{\addvspace{1pt}\sffamily\small}
{\contentslabel[\thecontentslabel]{1.25cm}}
{}
{\ \titlerule*[.5pc]{.}\;\thecontentspage}
[]

% Lista de figuras
\titlecontents{figure}[0em]
{\addvspace{-5pt}\sffamily}
{\figurename~\thecontentslabel.\hspace*{1em}}
{}
{\ \titlerule*[.5pc]{.}\;\thecontentspage}
[]

% Lista de tablas
\titlecontents{table}[0em]
{\addvspace{-5pt}\sffamily}
{\tablename~\thecontentslabel.\hspace*{1em}}
{}
{\ \titlerule*[.5pc]{.}\;\thecontentspage}
[]

%----------------------------------------------------------------------------------------
%	MINI TABLA DE CONTENIDO EN LAS PÁGINAS DE PARTE
%----------------------------------------------------------------------------------------

\titlecontents{lchapter}[0em]
{\addvspace{15pt}\large\sffamily\bfseries}
{\color{acento}\contentslabel[\Large\thecontentslabel]{1.25cm}\color{acento}}
{}
{\color{acento}\normalsize\sffamily\bfseries\;\titlerule*[.5pc]{.}\;\thecontentspage}

\titlecontents{lsection}[0em]
{\sffamily\small}
{\contentslabel[\thecontentslabel]{1.25cm}}
{}
{}

\titlecontents{lsubsection}[.5em]
{\normalfont\footnotesize\sffamily}
{}
{}
{}

%----------------------------------------------------------------------------------------
%	ENCABEZADOS DE PÁGINA
%----------------------------------------------------------------------------------------

\usepackage{fancyhdr}

\pagestyle{fancy}
\makeatletter
% En los preliminares (Resumen, Prólogo...) solo se muestra el título, sin "Capítulo 0"
\renewcommand{\chaptermark}[1]{\markboth{\sffamily\normalsize\bfseries\if@mainmatter\chaptername\ \thechapter.\ \fi#1}{}}
\makeatother
\renewcommand{\sectionmark}[1]{\markright{\sffamily\normalsize\thesection\hspace{5pt}#1}{}}
\fancyhf{} \fancyhead[LE,RO]{\sffamily\normalsize\thepage}
\fancyhead[LO]{\rightmark}
\fancyhead[RE]{\leftmark}
\renewcommand{\headrulewidth}{0.5pt}
\renewcommand{\headrule}{{\color{principal}\hrule width\headwidth height\headrulewidth\vskip-\headrulewidth}}
\renewcommand{\footrulewidth}{0pt}
\fancypagestyle{plain}{\fancyhead{}\renewcommand{\headrulewidth}{0pt}}

% Quita el encabezado de las páginas en blanco al final de los capítulos
\makeatletter
\renewcommand{\cleardoublepage}{
\clearpage\ifodd\c@page\else
\hbox{}
\vspace*{\fill}
\thispagestyle{empty}
\newpage
\fi}

%----------------------------------------------------------------------------------------
%	ESTILOS DE RECUADROS (definiciones, ejemplos, ejercicios...)
%----------------------------------------------------------------------------------------

\usepackage{amsmath,amsfonts,amssymb,amsthm}

\newtheorem{notation}{Notación}[chapter]

% Recuadros con encabezado de color
\newtheoremstyle{principalnumbox}
{0pt}{0pt}{\normalfont}{}
{\small\bf\sffamily\color{acento}}
{\;}{0.25em}
{\small\sffamily\color{acento}\thmname{#1}\nobreakspace\thmnumber{\@ifnotempty{#1}{}\@upn{#2}}%
\thmnote{\nobreakspace\the\thm@notefont\sffamily\bfseries\color{black}---\nobreakspace#3.}}
\renewcommand{\qedsymbol}{$\blacksquare$}

\newtheoremstyle{blacknumex}
{5pt}{5pt}{\normalfont}{}
{\small\bf\sffamily}
{\;}{0.25em}
{\small\sffamily{\tiny\ensuremath{\blacksquare}}\nobreakspace\thmname{#1}\nobreakspace\thmnumber{\@ifnotempty{#1}{}\@upn{#2}}%
\thmnote{\nobreakspace\the\thm@notefont\sffamily\bfseries---\nobreakspace#3.}}

\newtheoremstyle{blacknumbox}
{0pt}{0pt}{\normalfont}{}
{\small\bf\sffamily}
{\;}{0.25em}
{\small\sffamily\thmname{#1}\nobreakspace\thmnumber{\@ifnotempty{#1}{}\@upn{#2}}%
\thmnote{\nobreakspace\the\thm@notefont\sffamily\bfseries---\nobreakspace#3.}}

\newtheoremstyle{principalnum}
{5pt}{5pt}{\normalfont}{}
{\small\bf\sffamily\color{acento}}
{\;}{0.25em}
{\small\sffamily\color{acento}\thmname{#1}\nobreakspace\thmnumber{\@ifnotempty{#1}{}\@upn{#2}}%
\thmnote{\nobreakspace\the\thm@notefont\sffamily\bfseries\color{black}---\nobreakspace#3.}}
\makeatother

\newcounter{dummy}
\numberwithin{dummy}{section}
\theoremstyle{principalnumbox}
\newtheorem{theoremeT}[dummy]{Teorema}
\newtheorem{problem}{Problema}[chapter]
\newtheorem{exerciseT}{Ejercicio}[chapter]
\theoremstyle{blacknumex}
\newtheorem{exampleT}{Ejemplo}[chapter]
\theoremstyle{blacknumbox}
\newtheorem{vocabulary}{Vocabulario}[chapter]
\newtheorem{definitionT}{Definición}[section]
\newtheorem{corollaryT}[dummy]{Corolario}
\theoremstyle{principalnum}
\newtheorem{proposition}[dummy]{Proposición}

%----------------------------------------------------------------------------------------
%	RECUADROS DE COLOR
%----------------------------------------------------------------------------------------

\RequirePackage[framemethod=default]{mdframed}

% Teorema
\newmdenv[skipabove=7pt,skipbelow=7pt,backgroundcolor=black!5,linecolor=principal,
innerleftmargin=5pt,innerrightmargin=5pt,innertopmargin=5pt,innerbottommargin=5pt,
leftmargin=0cm,rightmargin=0cm]{tBox}

% Ejercicio
\newmdenv[skipabove=7pt,skipbelow=7pt,rightline=false,leftline=true,topline=false,bottomline=false,
backgroundcolor=principal!10,linecolor=principal,
innerleftmargin=5pt,innerrightmargin=5pt,innertopmargin=5pt,innerbottommargin=5pt,
leftmargin=0cm,rightmargin=0cm,linewidth=4pt]{eBox}

% Definición
\newmdenv[skipabove=7pt,skipbelow=7pt,rightline=false,leftline=true,topline=false,bottomline=false,
linecolor=principal,innerleftmargin=5pt,innerrightmargin=5pt,innertopmargin=0pt,innerbottommargin=0pt,
leftmargin=0cm,rightmargin=0cm,linewidth=4pt]{dBox}

% Corolario
\newmdenv[skipabove=7pt,skipbelow=7pt,rightline=false,leftline=true,topline=false,bottomline=false,
linecolor=gray,backgroundcolor=black!5,
innerleftmargin=5pt,innerrightmargin=5pt,innertopmargin=5pt,innerbottommargin=5pt,
leftmargin=0cm,rightmargin=0cm,linewidth=4pt]{cBox}

% Recuadro de presentación (texto destacado en los preliminares)
\newmdenv[skipabove=7pt,skipbelow=7pt,rightline=false,leftline=true,topline=false,bottomline=false,
backgroundcolor=principal!8,linecolor=principal,
innerleftmargin=12pt,innerrightmargin=12pt,innertopmargin=12pt,innerbottommargin=12pt,
leftmargin=0cm,rightmargin=0cm,linewidth=4pt]{destacado}

\newenvironment{theorem}{\begin{tBox}\begin{theoremeT}}{\end{theoremeT}\end{tBox}}
\newenvironment{exercise}{\begin{eBox}\begin{exerciseT}}{\hfill{\color{principal}\tiny\ensuremath{\blacksquare}}\end{exerciseT}\end{eBox}}
\newenvironment{definition}{\begin{dBox}\begin{definitionT}}{\end{definitionT}\end{dBox}}
\newenvironment{example}{\begin{exampleT}}{\hfill{\tiny\ensuremath{\blacksquare}}\end{exampleT}}
\newenvironment{corollary}{\begin{cBox}\begin{corollaryT}}{\end{corollaryT}\end{cBox}}

%----------------------------------------------------------------------------------------
%	NOTA AL MARGEN (círculo con "N")
%----------------------------------------------------------------------------------------

\newenvironment{remark}{\par\vspace{10pt}\small
\begin{list}{}{\leftmargin=35pt\rightmargin=25pt}\item\ignorespaces
\makebox[-2.5pt]{\begin{tikzpicture}[overlay]
\node[draw=principal!60,line width=1pt,circle,fill=principal!25,font=\sffamily\bfseries,inner sep=2pt,outer sep=0pt] at (-15pt,0pt){\textcolor{acento}{N}};\end{tikzpicture}}
\advance\baselineskip -1pt}{\end{list}\vskip5pt}

%----------------------------------------------------------------------------------------
%	TÍTULOS DE SECCIÓN (niveles APA 2 y 3, número en el margen)
%----------------------------------------------------------------------------------------

\makeatletter
\renewcommand{\@seccntformat}[1]{\llap{\textcolor{principal}{\csname the#1\endcsname}\hspace{1em}}}
% Nivel 2 APA: alineado a la izquierda, negrita
\renewcommand{\section}{\@startsection{section}{1}{\z@}
{-4ex \@plus -1ex \@minus -.4ex}
{1ex \@plus.2ex }
{\normalfont\large\sffamily\bfseries}}
% Nivel 3 APA: alineado a la izquierda, negrita cursiva
\renewcommand{\subsection}{\@startsection {subsection}{2}{\z@}
{-3ex \@plus -0.1ex \@minus -.4ex}
{0.5ex \@plus.2ex }
{\normalfont\sffamily\bfseries\itshape}}
\renewcommand{\subsubsection}{\@startsection {subsubsection}{3}{\z@}
{-2ex \@plus -0.1ex \@minus -.2ex}
{.2ex \@plus.2ex }
{\normalfont\small\sffamily\bfseries}}
\renewcommand\paragraph{\@startsection{paragraph}{4}{\z@}
{-2ex \@plus-.2ex \@minus .2ex}
{.1ex}
{\normalfont\small\sffamily\bfseries}}

%----------------------------------------------------------------------------------------
%	PÁGINAS DE PARTE (opcionales)
%----------------------------------------------------------------------------------------

\newcommand{\@mypartnumtocformat}[2]{%
\setlength\fboxsep{0pt}%
\noindent\colorbox{principal!20}{\strut\parbox[c][.7cm]{\ecart}{\color{acento}\Large\sffamily\bfseries\centering#1}}\hskip\esp\colorbox{principal!40}{\strut\parbox[c][.7cm]{\linewidth-\ecart-\esp}{\Large\sffamily\centering#2}}}%
\newcommand{\@myparttocformat}[1]{%
\setlength\fboxsep{0pt}%
\noindent\colorbox{principal!40}{\strut\parbox[c][.7cm]{\linewidth}{\Large\sffamily\centering#1}}}%
\newlength\esp
\setlength\esp{4pt}
\newlength\ecart
\setlength\ecart{1.2cm-\esp}
\newcommand{\thepartimage}{}%
\newcommand{\partimage}[1]{\renewcommand{\thepartimage}{#1}}%
\def\@part[#1]#2{%
\ifnum \c@secnumdepth >-2\relax%
\refstepcounter{part}%
\addcontentsline{toc}{part}{\texorpdfstring{\protect\@mypartnumtocformat{\thepart}{#1}}{\partname~\thepart\ ---\ #1}}
\else%
\addcontentsline{toc}{part}{\texorpdfstring{\protect\@myparttocformat{#1}}{#1}}%
\fi%
\startcontents%
\markboth{}{}%
{\thispagestyle{empty}%
\begin{tikzpicture}[remember picture,overlay]%
\node at (current page.north west){\begin{tikzpicture}[remember picture,overlay]%
\fill[principal!20](0cm,0cm) rectangle (\paperwidth,-\paperheight);
\node[anchor=north] at (4cm,-3.25cm){\color{principal!40}\fontsize{220}{100}\sffamily\bfseries\@Roman\c@part};
\node[anchor=south east] at (\paperwidth-1cm,-\paperheight+1cm){\parbox[t][][t]{8.5cm}{
\printcontents{l}{0}{\setcounter{tocdepth}{1}}%
}};
\node[anchor=north east] at (\paperwidth-1.5cm,-3.25cm){\parbox[t][][t]{15cm}{\strut\raggedleft\color{acento}\fontsize{30}{30}\sffamily\bfseries#2}};
\end{tikzpicture}};
\end{tikzpicture}}%
\@endpart}
\def\@spart#1{%
\startcontents%
\phantomsection
{\thispagestyle{empty}%
\begin{tikzpicture}[remember picture,overlay]%
\node at (current page.north west){\begin{tikzpicture}[remember picture,overlay]%
\fill[principal!20](0cm,0cm) rectangle (\paperwidth,-\paperheight);
\node[anchor=north east] at (\paperwidth-1.5cm,-3.25cm){\parbox[t][][t]{15cm}{\strut\raggedleft\color{acento}\fontsize{30}{30}\sffamily\bfseries#1}};
\end{tikzpicture}};
\end{tikzpicture}}
\addcontentsline{toc}{part}{\texorpdfstring{%
\setlength\fboxsep{0pt}%
\noindent\protect\colorbox{principal!40}{\strut\protect\parbox[c][.7cm]{\linewidth}{\Large\sffamily\protect\centering #1\quad\mbox{}}}}{#1}}%
\@endpart}
\def\@endpart{\vfil\newpage
\if@twoside
\if@openright
\null
\thispagestyle{empty}%
\newpage
\fi
\fi
\if@tempswa
\twocolumn
\fi}

%----------------------------------------------------------------------------------------
%	TÍTULOS DE CAPÍTULO (nivel 1 APA, con imagen de cabecera)
%----------------------------------------------------------------------------------------

\newif\ifusechapterimage
\usechapterimagetrue
\newcommand{\thechapterimage}{}%
\newcommand{\chapterimage}[1]{\ifusechapterimage\renewcommand{\thechapterimage}{#1}\fi}%
\def\@makechapterhead#1{%
{\parindent \z@ \raggedright \normalfont
\ifnum \c@secnumdepth >\m@ne
\if@mainmatter
\begin{tikzpicture}[remember picture,overlay]
\node at (current page.north west)
{\begin{tikzpicture}[remember picture,overlay]
\node[anchor=north west,inner sep=0pt] at (0,0) {\ifusechapterimage\includegraphics[width=\paperwidth]{\thechapterimage}\fi};
\draw[anchor=west] (\Gm@lmargin,-9cm) node [line width=2pt,rounded corners=15pt,draw=principal,fill=white,fill opacity=0.6,inner sep=15pt]{\strut\makebox[22cm]{}};
\draw[anchor=west] (\Gm@lmargin+.3cm,-9cm) node {\huge\sffamily\bfseries\color{black}\thechapter. #1\strut};
\end{tikzpicture}};
\end{tikzpicture}
\else
\begin{tikzpicture}[remember picture,overlay]
\node at (current page.north west)
{\begin{tikzpicture}[remember picture,overlay]
\node[anchor=north west,inner sep=0pt] at (0,0) {\ifusechapterimage\includegraphics[width=\paperwidth]{\thechapterimage}\fi};
\draw[anchor=west] (\Gm@lmargin,-9cm) node [line width=2pt,rounded corners=15pt,draw=principal,fill=white,fill opacity=0.6,inner sep=15pt]{\strut\makebox[22cm]{}};
\draw[anchor=west] (\Gm@lmargin+.3cm,-9cm) node {\huge\sffamily\bfseries\color{black}#1\strut};
\end{tikzpicture}};
\end{tikzpicture}
\fi\fi\par\vspace*{270\p@}}}

\def\@makeschapterhead#1{%
\begin{tikzpicture}[remember picture,overlay]
\node at (current page.north west)
{\begin{tikzpicture}[remember picture,overlay]
\node[anchor=north west,inner sep=0pt] at (0,0) {\ifusechapterimage\includegraphics[width=\paperwidth]{\thechapterimage}\fi};
\draw[anchor=west] (\Gm@lmargin,-9cm) node [line width=2pt,rounded corners=15pt,draw=principal,fill=white,fill opacity=0.6,inner sep=15pt]{\strut\makebox[22cm]{}};
\draw[anchor=west] (\Gm@lmargin+.3cm,-9cm) node {\huge\sffamily\bfseries\color{black}#1\strut};
\end{tikzpicture}};
\end{tikzpicture}
\par\vspace*{270\p@}}
\makeatother

%----------------------------------------------------------------------------------------
%	AUTORES (portadilla)
%----------------------------------------------------------------------------------------

% \autor{Nombre completo}{Afiliación, ORCID, ciudad, correo}
\newcommand{\autor}[2]{%
  \par\noindent{\sffamily\bfseries\large\color{acento}#1}\par
  {\small\begin{spacing}{1}\noindent #2\end{spacing}}\par\medskip}

%----------------------------------------------------------------------------------------
%	ENLACES DEL PDF
%----------------------------------------------------------------------------------------

\usepackage{hyperref}
\hypersetup{hidelinks,breaklinks=true,bookmarksopen=false}
\usepackage{bookmark}
\bookmarksetup{
open,
numbered,
addtohook={%
\ifnum\bookmarkget{level}=0 % capítulo
\bookmarksetup{bold}%
\fi
\ifnum\bookmarkget{level}=-1 % parte
\bookmarksetup{color=principal,bold}%
\fi
}
}
